\section{Choices to declare and limits}\label{sec:parameters}

The starting values and those chosen on the data are in Table~\ref{tab:notation}. The thesis
should state these choices, each argued in the section given:
\begin{itemize}
  \item $\varepsilon_z$ encodes the notion of ``together'': shared thermals if tight, visual
        contact if wide (Section~\ref{sec:delta});
  \item $\delta$ is bounded by the visual range through $v_\xc$ (Section~\ref{sec:delta}) and
        measured on the data at 2~min (Section~\ref{sec:delta-choice});
  \item $q=0.1$: at most one contact in ten is chance (Section~\ref{sec:q});
  \item $\varepsilon_h$, $\varepsilon_z$ and $g$ show no plateau on the data, so their values are
        definitions (Sections~\ref{sec:tolerances} and~\ref{sec:plateau});
  \item headings are compared only when both pilots glide; climb, search and unclassified fixes
        turn the condition off (Section~\ref{sec:roles}); \fix{in the run reported here the phases
        are Vilpellet's, on 1~Hz flights only, so the condition acts only between two such
        flights (Section~\ref{sec:data});}
  \item contacts before the departure are ignored: a group formed on the takeoff is recognised
        only after the exit from $R_\dep$ (Section~\ref{sec:local});
  \item grid positions are interpolated across at most $h_{\max}$, so tracks logged more coarsely
        enter with gaps (Section~\ref{sec:phase0});
  \item the signed median delay leaves out the instants of search, where the delay is not
        determined (Section~\ref{sec:lag});
  \item chance is defined by the null model, whose shifts of about an hour also remove the
        synchrony of pilots with the weather of the hour: what it calls real includes pilots who
        share the same thermals at the same time without following each other, the ``same air''
        notion of Remark~\ref{rem:notions}, which is not proof of interaction. $\tau^*$ is 23~min
        with shifts over an hour and 53~min with shifts under 30~min
        (Section~\ref{sec:shift-classes});
  \item the errors of the phase labelling reach the heading condition: a climb labelled as glide
        compares headings that turn, and joining the short gaps mitigates the problem.
\end{itemize}

The measurements have three further limits:
\begin{itemize}
  \item the profiles of Sections~\ref{sec:delta-choice} and~\ref{sec:tolerances} are on one sample
        of 300 dates;
  \item the plateau check of Section~\ref{sec:plateau} is on 300 dates and with the starting delay
        of 5~min;
  \item the data are the cleaned tracks of the old repository (soaring, tag
        \texttt{legacy-final}), not the pre-processing of xc-dynamics that
        Section~\ref{sec:data} describes.
\end{itemize}

\subsection{Checks done}\label{sec:checks-done}

\begin{itemize}
  \item \textbf{Code against this note.} Every phase was compared with its code (grid, departure,
        candidate filter, scan of delays, contact, gaps, signed delay, null model, configurations,
        episodes, membership, blocks): all implemented as written, with four differences of text,
        now corrected in red: (1) the phases of the run are Vilpellet's, on 1~Hz flights, so the
        heading condition acts only between two such flights; (2) headings are in the east and
        north of the first fix of each flight, not of the fix itself; (3) a contact is kept only
        from the later departure to the earlier landing; (4) instants in a gap of the track are
        labelled solo in the tables.
  \item \textbf{Properties of the tables} of the chosen run: every pair interval lasts at least
        $\tau^*=2285$~s and lies where both flights fly after departure; every configuration has at
        least two members; membership covers each flight from takeoff to landing without holes; on
        3\,000 flights the group time equals the union of its pair intervals; on 3\,000
        configurations the members are connected at the midpoint and no active interval leaves; on
        2\,000 episodes every configuration but the first starts when another ends and shares a
        member; every episode lies in one block. The counts of Table~\ref{tab:q} (154\,598 and
        6\,039 flights; 21.7\% and 6.9\% in a group) are reproduced.
\end{itemize}

\subsection{Other checks}\label{sec:checks}

\begin{itemize}
  \item Visual inspection: about ten episodes on a map, with the entries and exits checked by hand.
  \item Sensitivity to the step: on a subset of pairs, the computation with $\Delta=1$~s must give
        the same intervals.
  \item Competition days: if the dataset contains competition tracks, on those days everybody flies
        the same course and is in a group by construction; they are flagged and treated apart.
  \item Grid: the share of grid instants without a position, per flight and by logging interval.
\end{itemize}
